\chapter{Аналитическая часть}



В данном разделе рассматриваются рекурсивный и нерекурсивный подходы к решению задач, основанных на последовательной обработке числовых диапазонов. Рекурсивные алгоритмы реализуют принцип декомпозиции --- разбиение задачи на подзадачи того же типа, что обеспечивает простоту и математическую изящность описания. Нерекурсивные (итеративные) алгоритмы, напротив, опираются на циклы и явное управление состоянием, что повышает эффективность и снижает затраты памяти.



Рекурсивный алгоритм основан на принципе декомпозиции задачи, где основная проблема разбивается на более простые подзадачи того же типа. В данном случае алгоритм последовательно выводит текущее число и вызывает себя для следующего числа, пока не достигнет конечного значения. Такой подход демонстрирует наглядность математического решения, но имеет ограничения по глубине рекурсии~\cite{goloveshkin2020}.

Нерекурсивный алгоритм использует итеративный подход с явным управлением состоянием через цикл. Этот метод характеризуется последовательным перебором всех чисел в заданном диапазоне с минимальными накладными расходами.

Ключевое различие:

\begin{itemize}
	\item потребление памяти. Рекурсивный метод требует дополнительной памяти для хранения контекста вызовов в стеке, что может привести к его переполнению при больших диапазонах. Итеративный подход использует фиксированное количество памяти независимо от размера задачи~\cite{valedinsky2018}.
\end{itemize}


\section*{Вывод}

В аналитической части были рассмотрены рекурсивный и нерекурсивный алгоритмы, их принципы работы, особенности и различия в потреблении памяти.



\clearpage
